\subsection{Tools and Technologies}
\label{sec:tools_technologies}

Model development and training were implemented in Python using PyTorch
as the primary deep learning framework. PyTorch provides a dynamic computation
graph with imperative execution, enabling straightforward debugging and rapid
iteration during experimental design of oriented-detection training
loops~\cite{Paszke2019PyTorch}. Its native autograd engine supports
the gradient computation required for angular regression losses (e.g.,
Kullback--Leibler divergence loss over Gaussian-parameterised bounding boxes)
without requiring custom backward passes. Training experiments were executed
on Google Colaboratory accelerated hardware (NVIDIA T4 and A100
GPUs), a cloud-based environment that provides access to modern GPU architectures
without requiring local infrastructure investment (a constraint particularly
relevant for research groups in resource-limited institutions)~\cite{Garavagno2022ColabNAS}.

Model architectures are drawn from the Ultralytics YOLO ecosystem~\cite{Jocher2026}, which provides YOLO11\allowbreak-OBB and YOLO26\allowbreak-OBB with built-in support
for oriented bounding box regression, rotated-IoU-based matching, and
rotation-aware data augmentation. The Ultralytics library also exposes
an internal rotated-IoU computation that was reused as a primitive within
the custom Macro AP-rIoU@[0.50:0.80] evaluator implemented for this work.
For inference-speed benchmarking on edge\allowbreak-representative hardware, models were
exported to the ONNX portable format, enabling CPU-only evaluation
independent of the PyTorch runtime. Preprocessing utilities, annotation
parsing, clip-level dataset splitting, and result formatting were implemented
as auxiliary Python scripts to ensure full reproducibility across experimental
configurations.
